% This file is intended to be included via \input{} or \include{}.
% It contains no preamble and no document environment.

\section*{Math Primer for Appendix A}

 Math Primer for Understanding Appendix A\par

This primer introduces the minimum mathematical concepts needed to understand Appendix A of the FRFP paper.
Everything is presented in poster-style blocks, with motivation first, intuition second, and formal meaning last.\par

No prior category theory is assumed.\par

\section*{1. OBJECTS}

\subsection*{[PHASE] BLOCK 1 --- What Is an Object?}

Motivation: We need to talk about things before we can talk about transforming them.\par

An object is simply a thing we care about.\par

\n\n\begin{itemize}\n\item a number
\item a vector
\item a workflow
\item a pipeline\n\end{itemize}

Note: Key Idea: Objects are states, configurations, or entities --- not actions.\par

\subsection*{[PHASE] BLOCK 2 --- Objects in FRFP}

Motivation: The AI must reason about explicit artifacts at multiple levels of granularity.\par

In FRFP, there are two related but distinct senses in which we talk about objects:\par

Atomic objects: RI, EC, ED\par

Composite objects: explicit pipelines built from RI, EC, ED\par

Both are objects, but they play different roles.\par

Every pipeline is a thing the system can:\n\n\begin{itemize}\n\item inspect
\item transform
\item compare\n\end{itemize}

\section*{2. MORPHISMS (TRANSFORMATIONS)}

\subsection*{[PHASE] BLOCK 1 --- What Is a Morphism?}

Motivation: Reasoning is not static --- things change.\par

A morphism is an allowed transformation from one object to another.\par

Written as:\par

A $\to$ B\par

Meaning: > There is a permitted way to go from A to B. \par

\subsection*{[PHASE] BLOCK 2 --- Morphisms as Reasoning Steps}

Motivation: Reasoning is not static --- but not everything is transformable in the same way.\par

In FRFP, morphisms primarily act on composite objects (pipelines), not on atomic objects directly.\par

Pipelines change via explicit reductions\par

Atomic objects serve as stable building blocks\par

Note: Key Idea: Objects exist at multiple levels, but not all levels have interesting morphisms.\par

Formal clarification: Atomic objects (RI, EC, ED) admit only identity morphisms, while all nontrivial morphisms of the explicit category E act on composite objects (explicit pipelines) built from them.\par

Motivation: To model reasoning, we must model steps.\par

In FRFP:\n\n\begin{itemize}\n\item morphisms represent explicit reasoning steps
\item each step is inspectable and rule-governed\n\end{itemize}

Note: Objects = states of reasoning
Note: Morphisms = steps of reasoning\par

\section*{3. COMPOSITION}

\subsection*{[PHASE] BLOCK 1 --- Why Composition Matters}

Motivation: Reasoning happens in multiple steps.\par

If:\par

A $\to$ B
B $\to$ C\par

then we can compose them:\par

A $\to$ C\par

\subsection*{[PHASE] BLOCK 2 --- Composition Is Associative}

Motivation: Grouping steps should not affect outcomes.\par

Doing:\par

(A $\to$ B) $\to$ C\par

is the same as:\par

A $\to$ (B $\to$ C)\par

 This property makes long chains of reasoning manageable.\par

\section*{4. IDENTITY MORPHISMS (STRUCTURAL IDENTITIES)}

\subsection*{[PHASE] BLOCK 1 --- Doing Nothing Is Still a Step}

Motivation: Any formal system must allow no change.\par

For every object A, there is an identity morphism:\par

id\_A : A $\to$ A\par

Meaning: > Leave A exactly as it is. \par

\subsection*{[PHASE] BLOCK 2 --- Why Structural Identities Matter}

Motivation: Without identities, composition breaks.\par

Identities guarantee:\n\n\begin{itemize}\n\item pipelines can remain unchanged
\item composition laws hold
\item equivalence can be defined\n\end{itemize}

Note: In FRFP: Structural identities mean the same pipeline is still the same pipeline.\par

\section*{5. ALGEBRAIC STRUCTURE}

\subsection*{[PHASE] BLOCK 1 --- What Does Algebraic Mean Here?}

Motivation: We want to prove things about reasoning, not just run it.\par

 Algebraic structure means:\n\n\begin{itemize}\n\item objects
\item transformations
\item rules for combining transformations
\item predictable behavior\n\end{itemize}

It does not mean equations or numbers.\par

\subsection*{[PHASE] BLOCK 2 --- Why Algebraic Structure Enables Formal Reasoning}

Motivation: Proofs require stable rules.\par

Because the system has algebraic structure, we can prove:\n\n\begin{itemize}\n\item termination
\item normalization
\item equivalence
\item correctness preservation\n\end{itemize}

Key point: Without algebraic structure, none of these claims are meaningful.\par

\section*{6. CATEGORIES (THE MINIMAL FRAMEWORK)}

\subsection*{[PHASE] BLOCK 1 --- What Is a Category?}

Motivation: We want the smallest structure that supports formal reasoning.\par

A category consists of: 1. objects 2. morphisms between objects 3. composition of morphisms 4. identity morphisms\par

That s it.\par

\subsection*{[PHASE] BLOCK 2 --- Why Categories Appear in FRFP}

Motivation: FRFP needs structure, not decoration.\par

Categories allow us to:\n\n\begin{itemize}\n\item model workflows as objects
\item model reasoning as morphisms
\item guarantee compositional correctness\n\end{itemize}

 Category theory is used because it is the weakest tool that works.\par

\section*{7. GROUPS (FOR CONTRAST)}

\subsection*{[PHASE] BLOCK 1 --- What Is a Group?}

Motivation: To understand categories and navigation, it helps to see what they generalize.\par

A group has:\n\n\begin{itemize}\n\item a single underlying object
\item elements acting as transformations of that object
\item an identity element
\item inverses for every element\n\end{itemize}

Everything is reversible and acts on one thing.\par

Note: Example intuition: Rotations of a square, permutations of a fixed set.\par

\subsection*{[PHASE] BLOCK 2 --- Why FRFP Does Not Use Groups}

Motivation: Reasoning and restructuring do not act on a single fixed object.\par

Explicit reduction:\n\n\begin{itemize}\n\item simplifies
\item loses syntactic detail
\item is directional\n\end{itemize}

Navigation:\n\n\begin{itemize}\n\item is invertible
\item but moves between different shapes\n\end{itemize}

Note: Hence: FRFP needs something more general than groups.\par

\section*{8. GROUPOIDS (THE RIGHT NOTION FOR NAVIGATION)}

\subsection*{[PHASE] BLOCK 1 --- What Is a Groupoid?}

Motivation: We need reversibility without collapsing everything to one object.\par

A groupoid is:\n\n\begin{itemize}\n\item a category
\item with many objects
\item where every morphism is invertible\n\end{itemize}

Note: Intuition: A groupoid is a group with many objects. \par

\subsection*{[PHASE] BLOCK 2 --- Why Navigation Forms a Groupoid}

Motivation: Structural refactorings must be reversible but may change form.\par

In FRFP navigation:\n\n\begin{itemize}\n\item objects = shapes
\item morphisms = invertible structural rewrites
\item rewrites go between different shapes\n\end{itemize}

Each rewrite has an inverse, but there is no single global shape.\par

Note: Conclusion: Navigation is a groupoid, not a group.\par

\subsection*{[PHASE] BLOCK 3 --- One-Sentence Formal Clarification}

Navigation is modeled as a groupoid rather than a group because structural rewrites are invertible but act between many distinct shapes, not as symmetries of a single shape.\par

\section*{9. WHY CARTESIAN LIFTING IS A CORE HUMAN--AI SAFETY IDEA}

\subsection*{[PHASE] BLOCK 1 --- Why Lifting Is Needed at All}

Motivation: If structure and content are independent, changing structure should not change computation --- so why do anything?\par

The subtlety is this:\n\n\begin{itemize}\n\item structure and content are independent in meaning
\item but content is still located inside structure\n\end{itemize}

When structure changes, the meaning must stay the same, but the placement of content must be updated.\par

\subsection*{[PHASE] BLOCK 2 --- What Goes Wrong If Lifting Is Implicit}

Motivation: Hidden mechanisms break trust.\par

If lifting were implicit:\n\n\begin{itemize}\n\item the AI could reinterpret content during restructuring
\item structure changes could smuggle in computation
\item meaning preservation would be an assumption, not a theorem\n\end{itemize}

 Silent movement is silent reasoning.\par

\subsection*{[PHASE] BLOCK 3 --- What Explicit Lifting Guarantees}

Motivation: Humans need refactoring to be as safe as mathematics.\par

By making lifting:\n\n\begin{itemize}\n\item explicit (it is a visible operation)
\item canonical (there is only one correct way)
\item unique (no ambiguity)\n\end{itemize}

you get:\n\n\begin{itemize}\n\item semantic invariance (meaning does not change)
\item confluence (order of edits does not matter)
\item human trust (nothing hidden can happen)\n\end{itemize}

 Key Insight: Cartesian lifting plays the same safety role for structure/content separation that the tacit/explicit boundary plays for human/AI separation.\par

\subsection*{[PHASE] BLOCK 4 --- One-Sentence Takeaway}

Cartesian lifting exists because structure changes the placement of computation even when it must not change its meaning, and lifting is the unique, explicit rule that updates this interface safely.\par

\section*{SUPER SUMMARY --- HOW THIS PRIMER CONNECTS TO APPENDIX A}

Objects $\to$ explicit objects (atomic RI/EC/ED and composite pipelines)\par

Morphisms $\to$ explicit reasoning steps\par

Composition $\to$ multi-step reasoning\par

Identity $\to$ structural sameness\par

Algebraic structure $\to$ provable properties\par

Category $\to$ minimal framework that supports all of the above\par

 This primer provides exactly the tools needed to understand Appendix A --- nothing more, nothing less.\par

\section*{HOW TO READ SYMBOLS (A POSTER-STYLE GUIDE)}

This page explains how to read the symbols used throughout Appendix A and the math primer. Nothing new is introduced --- this is purely about interpretation.\par

\subsection*{[PHASE] BLOCK 1 --- Objects and Membership}

Symbol:\par

e E\par

Read as:\par

 e is an explicit pipeline (an object in the explicit category). \par

Informal meaning:\par

e is a concrete, inspectable workflow\par

it lives entirely in the explicit world\par

\subsection*{[PHASE] BLOCK 2 --- Morphisms (Arrows)}

Symbol:\par

e $\to$ e'\par

Read as:\par

 There is an allowed explicit transformation from pipeline e to pipeline e'. \par

Informal meaning:\par

one step (or many steps) of explicit reasoning\par

often a simplification or normalization\par

\subsection*{[PHASE] BLOCK 3 --- Explicit Reduction}

Symbol:\par

e $\to\_\{e\}$ e'\par

Read as:\par

 Pipeline e explicitly reduces to pipeline e'. \par

Informal meaning:\par

computation happened\par

expressions simplified\par

no structural reorganization yet\par

\subsection*{[PHASE] BLOCK 4 --- Identity Morphisms}

Symbol:\par

id\_e : e $\to$ e\par

Read as:\par

 The pipeline e stays exactly the same. \par

Informal meaning:\par

a formal way of saying do nothing \par

required so pipelines behave like mathematical objects\par

\subsection*{[PHASE] BLOCK 5 --- Composition of Morphisms}

Symbol:\par

e $\to$ e' $\to$ e \par

Read as:\par

 First transform e into e', then transform e' into e . \par

Informal meaning:\par

multi-step reasoning\par

chaining explicit steps\par

\subsection*{[PHASE] BLOCK 6 --- Categories}

Symbol:\par

E, S, N\par

Read as:\par

 A collection of objects together with allowed transformations between them. \par

Informal meaning in FRFP:\par

E = explicit pipelines + explicit reductions\par

S = shapes + structure-preserving maps\par

N = navigation steps (invertible structural rewrites)\par

\subsection*{[PHASE] BLOCK 7 --- Fibers}

Symbol:\par

E\_S = \{ e E | shape(e) = S \}\par

Read as:\par

 All explicit pipelines whose structure is S. \par

Informal meaning:\par

a folder (S)\par

containing all pipelines with that structure\par

\subsection*{[PHASE] BLOCK 8 --- Functions and Functors}

Symbol:\par

p : E $\to$ S\par

Read as:\par

 A rule that maps each explicit pipeline to its shape. \par

Informal meaning:\par

a consistent translation from one kind of object to another\par

preserves relationships between objects\par

\subsection*{[PHASE] BLOCK 9 --- Semantic Interpretation}

Symbol:\par

\(\llbracket$e\rrbracket\)\par

Read as:\par

 The human-interpretable meaning of pipeline e. \par

Informal meaning:\par

what the explicit computation means\par

evaluated by humans using tacit judgment\par

\subsection*{[PHASE] BLOCK 10 --- Navigation Action}

Symbol:\par

$\sigma \cdot e$ = e'\par

Read as:\par

 Apply the structural rewrite $\sigma$ to pipeline e, producing pipeline e'. \par

Informal meaning:\par

move the workflow to a new structure\par

content follows structure safely\par

[PHASE] SUPER SUMMARY --- HOW TO READ THE MATH\par

Symbols describe:\par

what exists (objects)\par

what is allowed (morphisms)\par

how steps combine (composition)\par

how structure and content relate (functions, functors)\par

how meaning is assigned (\(\llbracket$ \rrbracket\))\par

 Rule of thumb: Every symbol in Appendix A refers to an explicit, inspectable object or transformation --- never to hidden reasoning.\par

End of How to Read Symbols Poster Page\par

\section*{CONCEPT MAP --- HOW FRFP IS BUILT (DUAL-FLOW VIEW)}

This poster-style diagram shows two parallel conceptual flows --- content and structure --- and where they are formally synchronized. Read top to bottom.\par

CONTENT FLOW (EXPLICIT REASONING) STRUCTURE FLOW (WORKFLOW FORM)

 
 OBJECTS TDG / SHAPES 
 (What the AI reasons on) (Pure structure only) 
 
 Atomic: RI, EC, ED Grammar \& skeletons 
 Composite: Pipelines 
 
 
 
 
 MORPHISMS SHAPE CATEGORY S 
 (Explicit reductions $\to\_\{e\}$) (Shapes as objects) 
 
 Identity + computation Structure-preserving 
 
 
 
 
 COMPOSITION NAVIGATION N 
 (Chaining reductions) (Invertible rewrites) 
 
 e $\to$ e' $\to$ e Reorder / refactor form 
 
 
 
 
 ALGEBRAIC STRUCTURE 
 (Rule-governed behavior) 
 
 Normalization possible 
 
 
 
 
 CATEGORY E 
 (Explicit world) 
 
 Pipelines + $\to\_\{e\}$ 
 
 
 
 
 
 FIBRATION p : E $\to$ S 
 (Content structure) 
 
 Pipelines live over 
 shapes (fibers) 
 
 
 
 
 CARTESIAN LIFTING 
 (Safe synchronization) 
 
 Structure moves 
 content follows 
 
 
 
 
 COMBINED SYSTEM N $\ltimes$ E 
 (Structure + content) 
 
 Morphisms: ($\sigma$, f) 
 
 
 
 
 SEMANTICS 
 \(\llbracket$ \rrbracket\) , RB, TE, HFD 
 
 Explicit $\to$ Meaning 
 
 
 
 
 FRFP THEOREMS 
 (What we can prove) 
 
 \item Normalization 
 \item Confluence 
 \item Correctness 
 \item Semantic invariance 
 \par

[PHASE] HOW TO READ THIS DIAGRAM\par

The left column develops explicit content and computation\par

The right column develops pure structure and reorganization\par

They are intentionally separate until the fibration\par

The fibration is the first synchronization point\par

Everything below it depends on both flows\par

Key point: Key insight: Structure never changes meaning, and content never dictates structure --- they are coordinated only via principled mechanisms.\par

[PHASE] ONE-SENTENCE BIG PICTURE\par

FRFP develops content and structure independently, synchronizes them via fibration and lifting, and only then proves global correctness and semantic stability.\par

End of Concept Map Poster Page\par
